\lhead[ \leftmark ]{Informatik, Ergänzungsfach Informatik}
\rhead[Informatik]{Thomas Graf, \today, Winterthur}

\section*{Ergänzungsfach Informatik}

Im Ergänzungsfach Informatik (EF Informatik) werden fundierte Kenntnisse und praktische Fähigkeiten in essenziellen Bereichen der Computerwissenschaften vermittelt.
Wir verstehen das EF Informatik als eine höchst effektive Vorbereitung auf das Studium der Naturwissenschaften, der Ingenieurswissenschaften sowie der Mathematik und der Informatik.
Von grundlegenden Konzepten bis hin zu fortgeschrittenen Anwendungen bietet dieser Kurs eine ausgewogene Mischung aus Theorie und Praxis. Die Inhalte richten sich an Schülerinnen und Schüler, die ein starkes Interesse an Problemlösung haben.

\subsection*{Übersicht der Inhalte}
\begin{itemize}
	\item \textbf{praktische Werkzeuge der Informatik}\\ Textverarbeitung in \LaTeX\ zum Verfassen wissenschaftlicher Arbeiten (e.g. Maturitätsarbeit, Bachelorarbeit), Versionskontrolle und Kollaboration mithilfe von Git / GitHub
	\item \textbf{Datenstrukturen}\\ Einführung in effiziente Speicher- und Zugriffsmethoden wie Arrays, Listen, Bäume und Graphen. Schwerpunkt auf der Auswahl geeigneter Datenstrukturen zur Lösung spezifischer Probleme.
	\item \textbf{Rekursion}\\ Verständnis der Rekursion als mächtiges Programmierparadigma. Anwendungsbeispiele wie Suchalgorithmen, Sortierverfahren und rekursive Problemlösungsstrategien, Lösen von Sudokus.
	\item \textbf{Algorithmen}\\ Entwicklung und Analyse von Algorithmen
	\item \textbf{Komplexität}\\ Einführung in die Analyse der Laufzeit und Effizienz von Algorithmen, Definition der Komplexitätsklassen
	\item \textbf{Objektorientierte Programmierung (OOP)}\\ Prinzipien der OOP, wie Klassen, Vererbung, Polymorphismus und Abstraktion. Entwicklung modularer, wiederverwendbarer und wartbarer Programme. Wie wird ein grösseres Softwareprojekt aufgebaut?
	\item \textbf{Modellierung und Simulation}\\ Modellierung und Simulation dynamischer Systeme mithilfe des Verlet-Algorithmus, Einführung in die Techniken zur Erstellung von Animationsfilmen
	\item \textbf{Künstliche Intelligenz (KI)}\\
	      Einführung in die Prinzipien und Anwendungen von KI, einschliesslich maschinellem Lernen, neuronalen Netzen und Entscheidungsmodellen. Praktische Übungen mit Tools und Bibliotheken wie TensorFlow oder scikit-learn.
\end{itemize}
